\documentclass[aspectratio=169,11pt,usepdftitle=false]{beamer}
% Spectrolaminar column deck; storyboard main.md; figures from make_figs.py (records in the last page).
\usepackage[T1]{fontenc}
\usepackage[scaled=0.95]{helvet}
\renewcommand{\familydefault}{\sfdefault}
\usepackage{graphicx,booktabs,amsmath,amssymb}
\usefonttheme{professionalfonts}
\usepackage{sansmath}\sansmath
\graphicspath{{figs/}}
\hypersetup{pdftitle={Testing cell-type density as a circuit mechanism of primate spectrolaminar oscillations},pdfauthor={Hamed Nejat}}

\definecolor{gold}{HTML}{CFB87C}
\definecolor{pagebg}{HTML}{000000}
\definecolor{textfg}{HTML}{F2F2F2}
\definecolor{dim}{HTML}{A6A6A6}
\setbeamercolor{background canvas}{bg=pagebg}
\setbeamercolor{normal text}{fg=textfg,bg=pagebg}
\setbeamercolor{frametitle}{fg=gold,bg=pagebg}
\setbeamercolor{title}{fg=gold}
\setbeamercolor{structure}{fg=gold}
\setbeamerfont{frametitle}{size=\large,series=\bfseries}
\setbeamertemplate{navigation symbols}{}
\setbeamertemplate{frametitle}{\vspace{0.25em}\insertframetitle\par\vspace{-0.2em}}
\setbeamersize{text margin left=0.45cm,text margin right=0.45cm}
\setbeamertemplate{footline}{\hspace{0.45cm}{\color{dim}\fontsize{7}{8}\selectfont H. Nejat $\cdot$ Cell-type density
  and the spectrolaminar motif\hfill\insertframenumber/\inserttotalframenumber}\hspace{0.45cm}\vspace{0.18cm}}

% figure box: white card so white-ground figures read on black
\newcommand{\card}[2]{{\setlength{\fboxsep}{2pt}\colorbox{white}{\includegraphics[#1]{#2}}}}
\newcommand{\take}[1]{\par\vfill{\color{gold}\small #1}\par}
\newcommand{\src}[1]{{\color{dim}\scriptsize #1}\par\vspace{0.05cm}}

\begin{document}

% 1
\begin{frame}[plain]
\vfill
{\color{gold}\LARGE\bfseries Testing cell-type density as a circuit mechanism\\[0.2em]of primate spectrolaminar oscillations\par}
\vspace{1.0em}
{\normalsize A column with density-sized units and an imposed SOM/PV ramp puts alpha-beta in the deep half on the reference seeds, and leaves gamma near layer 4 weak. A three-source readout could carry the observed gamma profile; readout depth remains a confound.\par}
\vspace{1.4em}
{\normalsize Hamed Nejat\par}
{\color{dim}\small Bastos Lab $\cdot$ Vanderbilt University $\cdot$ November 10, 2026\par}
\vfill
{\color{dim}\small Preliminary, unpublished results.\par}
\vspace{0.4cm}
\end{frame}

% ---------- 1. The motif and how the model makes it ----------
% 2
\begin{frame}{Macaque cortex: gamma is relatively strongest above L4, alpha-beta below it}
\centering
\card{height=5.7cm}{densities.pdf}\hspace{0.2cm}\card{height=5.7cm}{obs_maps.pdf}
\take{Left: cell densities by depth (median of 11 areas). Right: LFP power / its maximum over depth.}
\src{LANSA data, Mulvey et al.\ 2024; motif after Mendoza-Halliday et al.\ 2024.}
\end{frame}

% 3
\begin{frame}{Unit sizes come from these densities; the SOM/PV ramp is imposed}
\centering
\card{height=5.15cm}{column.pdf}\hspace{0.25cm}\card{height=5.15cm}{built.pdf}
\take{4000 E cells per unit on average; 14\,251 cells (S 6870, M 4356, D 3025). SOM/PV 0.25, 1, 2: imposed.}
\src{AdEx units at $-0.4$, 0, $+0.3$; in-degree scales with density by an assumed rule. E/PV/SOM: S 5785/868/217, M 3668/344/344, D 2547/159/319. Built PV follows unit size $\times$ $\rho$; SOM follows the ramp; CB is calbindin, which labels more than SST cells. NG: pyramidal proxy.}
\end{frame}

% 4
\begin{frame}{Isolated, SOM/PV lowers the rhythm; in the column, M and D peak in alpha-beta}
\centering
\card{height=5.15cm}{mechanism.pdf}
\take{At the imposed ratios, M and D peak at 11--12 Hz in the column, S at 43 Hz (below the scored band); cause inferred: coupling and density scaling.}
\src{Matches its reference (34.0 vs $33.8\pm1$ Hz); $-6.27$ Hz/octave for $\rho$ 0.25--2; $\rho=4$ silent: inconclusive as specified.}
\end{frame}

% 5
\begin{frame}{The column fires sparsely: E cells below 1.6 Hz, interneurons faster}
\centering
\card{height=5.9cm}{ref_dynamics.png}
\take{Rates over 3 seeds: E 0.39--1.58 Hz, PV 5.6--22.6 Hz, SOM 4.7--9.4 Hz (lowest at S, highest at D).}
\src{Reference column, seed 1; 2 s of the 8 s run shown (first 2 s discarded). Left: sampled cells; right: per-unit LFP proxy (weighted synaptic currents).}
\end{frame}

% 6
\begin{frame}{The column puts alpha-beta in the deep half and most gamma at S}
\centering
\card{height=5.5cm}{ref_triplet.pdf}
\take{$r_{ab}$ 0.83--0.86 on seeds 1--3; flat density gives 0.84 as well.}
\src{$r_{ab}$, $r_{gm}$: depth correlation of the 10--30 / 50--150 Hz profile with the observed one. Model maps: LFP proxies read through the three-source Gaussian kernel, $\sigma=0.15$. Seed means, 3 seeds; stationary Poisson drive, whereas the observed maps are peri-stimulus. The reference was not rerun on fresh seeds; the later selected setting lost $r_{ab}$ there (0.35--0.47).}
\end{frame}

% 7
\begin{frame}{Near layer 4, gamma is weak: 0.19--0.31 in 12 runs against 0.848 observed}
\centering
\card{height=6.0cm}{mid_gamma.pdf}
\hspace{0.3cm}
\begin{minipage}[b]{0.33\textwidth}\raggedright\small
Cross terms carry 0--12\% of this power; M supplies 72--81\%.\\[0.6em]
S holds the maximum at 97--101 of 101 gamma frequencies.
\end{minipage}
\take{M's gamma is small relative to S's.}
\src{9 size-ladder runs ($N$ 1000/2000/4000 E per unit $\times$ 3 seeds) and 3 confirmation runs; 50--150 Hz, depths $-0.2$ to 0.}
\end{frame}

% ---------- 2. What disrupts it ----------
% 8
\begin{frame}{Pass rule, fixed before the parameter search: band and whole-map criteria}
\centering
\card{height=5.6cm}{scorer_triplet.pdf}
\take{Pass = crossing $>2$ SE, $r_{ab}, r_{gm}\ge0.6$, true unit order first, map $r$ above null and control, halves agree, result finite.}
\src{Crossing: the alpha-beta minus gamma profile is below $-2$ SE at the top contact and above $+2$ SE at the bottom. Halves: $|r_1 - r_2| \le 2\,\mathrm{SE}_\mathrm{diff}$, jackknife SEs over windows in each half. Null: random smooth unit spectra (95th percentile); control: every unit given S's spectrum. Shown: $B=1.41$, SOM$\to$E 15 ms in all units, seed 1: map $r$ 0.36, $r_{ab}$ 0.75, $r_{gm}$ 0.42; fails.}
\end{frame}

% 9
\begin{frame}{No run passes; $r_{gm}\ge0.6$ ($\sigma$ 0.15) only in 3 weight runs, none with the crossing}
\centering
\card{height=5.6cm}{value_scatter.pdf}
\hspace{0.3cm}
\begin{minipage}[b]{0.32\textwidth}\raggedright\small
87 scores, 78 distinct runs: size ladder 9, search 63 (54 new runs), confirmation 3, weight runs 12 (13 settings with their start, page 11). The 12 control runs are on pp.\ 12--13.\\[0.6em]
Largest $r_{gm}$ outside the weight runs: 0.57.\\[0.6em]
Whole-map ceiling, map $r$ ($r_\text{proto}$): 0.854 coherent, 0.693 independent.
\end{minipage}
\take{Gamma-profile ceiling of the readout: $r_{gm}$ 0.853 coherent, 0.657 independent ($\sigma=0.15$).}
\src{Ceiling = the best any source covariance reaches, from the observed data alone. At $\sigma$ 0.20 / 0.25, $r_{gm}\ge0.6$ in 9 and 11 of the 12 ladder and confirmation runs (none reaches the threshold at 0.15); weak mid-depth gamma holds at every width.}
\end{frame}

% 10
\begin{frame}{Tuning local in-degree trades alpha-beta against gamma; no level passes}
\centering
\card{width=13.4cm}{t1_kd_atlas.pdf}
\take{In-degree $\propto$ density$^B$; raising $B$ trades $r_{ab}$ for $r_{gm}$. $B$ 1.41 kept.}
\src{7 levels $\times$ 3 seeds, $N=4000$, 8 s. Selection on whole-map similarity ($r_\text{proto}$), not band criteria. Final selection ($B$ 1.41, SOM$\to$E 21 ms in D), fresh seeds: $r_{ab}$ 0.35--0.47, 0/3 pass.}
\end{frame}

% 11
\begin{frame}{Weight changes (mainly at M): $r_{gm}\ge0.6$ only without the crossing}
\centering
\card{height=5.2cm}{weights.pdf}
\take{$r_{gm}\ge0.6$ only without the crossing (A); with the crossing kept, $r_{gm}$ 0.42 $\to$ at most 0.54 (B; $r_{ab}$ 0.78). All 13 settings fail the scorer.}
\src{Exploratory, seed 1, 8 s; base $B$ 1.41, SOM$\to$E 15 ms in all units (page 8). A: M E$\to$E $\times$1.68, S, M E$\to$SOM $\times$0.5. B: M E$\to$E $\times$1.68, S E$\to$SOM $\times$0.5, M PV$\to$E $\times$1.25. M's recurrent excitation sets both bands (inferred).}
\end{frame}

% 12
\begin{frame}{Changing density or the imposed SOM/PV ramp breaks the frequency split}
\centering
\card{width=13.4cm}{t8_atlas.pdf}
\take{Density arms: D dominates most frequencies; SOM/PV arms: S does (read from the maps).}
\src{Seed means, 3 seeds; controls descriptive, uncorrected. One rule changed per control; total cells fixed. Flat ratio 0.79 = the ramp's geometric mean.}
\end{frame}

% 12
\begin{frame}{Each change hands the map to one unit; none raises gamma near L4}
\centering
\card{height=5.6cm}{t8_summary.pdf}
\take{No arm raises gamma near L4 (0.04--0.23; reference 0.23--0.26); flat density keeps $r_{ab}$ through deep dominance (inferred); the ratio arms' $r_{gm}$ near 0.5 comes from S dominating most frequencies.}
\src{Black: reference; colour: dominant unit (D cyan, S blue). Ratio = imposed ramp; arms also change sizes, in-degrees, rates.}
\end{frame}

% ---------- 3. Maximal similarity, step by step ----------
\begin{frame}{Is the readout the bottleneck?}
\vfill
\centering
\begin{minipage}{0.85\textwidth}\raggedright\large
No circuit setting, control or weight change passed the scorer or reached the observed mid-depth gamma (0.848).\\[1.2em]
Next, without simulation: could any mix of three depth sources reproduce it? If none can, the readout is the bottleneck at the declared depths. Whether those depths are right is tested last.
\end{minipage}
\vfill
\end{frame}

% 13
\begin{frame}{Step 1: target = observed gamma profile; readout = three Gaussian sources}
\centering
\begin{minipage}[c]{0.58\textwidth}\centering\card{height=5.0cm}{step1_target_kernel.pdf}\end{minipage}\hfill
\begin{minipage}[c]{0.39\textwidth}\raggedright
\[ P_z = k_z^{\top} C\, k_z,\quad C = LL^{\top}\succeq 0 \]
\[ k_{zu} = e^{-(z-z_u)^2/2\sigma^2} \]
\[ R_z = P_z / \max_z P_z \]
{\small $z$: contact depth (21); $u\in\{S,M,D\}$; $C$: source covariance, one for 50--150 Hz; mid = mean $R_z$, $z\in[-0.2,0]$ (gold band).}
\end{minipage}
\take{No simulation: search $C$ for the best fit this readout can represent.}
\src{Target $y$ = observed 50--150 Hz band mean / max (mid 0.946); per-frequency map mid 0.848.}
\end{frame}

% 14
\begin{frame}{Step 2: a single source at M caps mid-depth gamma at 0.615}
\centering
\card{height=6.0cm}{step2_profiles.pdf}
\take{$\sigma=0.15$. Raising M's power, even $10^4$-fold, approaches this kernel limit, 0.615, below the observed 0.848.}
\src{Curve: band-mean profile / its maximum (mid 0.946); bars: per-frequency map (mid 0.848); different normalisations.}
\end{frame}

% 15
\begin{frame}{Step 3: independent sources reach $r=0.657$, with mid-depth gamma 0.49}
\centering
\card{height=6.0cm}{step3_profiles.pdf}
\take{$\sigma=0.15$. Best diagonal $C$: S 58\%, M 42\%, D 0\%; $r_{gm}\ge0.6$ is passable without coherence.}
\src{48 optimizer starts.}
\end{frame}

% 16
\begin{frame}{Step 4: coherent sources reach $r=0.853$; that optimum has mid-depth 0.72}
\centering
\card{height=6.0cm}{step4_profiles.pdf}
\take{$\sigma=0.15$. Controls: 5 random profiles recovered at $r\ge0.99999$; alternating profile $r$ 0.12.}
\src{A least-squares fit gives only $r$ 0.45. Optimum rank one: S--M coherence $+1$, D $-1$.}
\end{frame}

% 17
\begin{frame}{Step 5: at $\sigma=0.15$ the observed mid-depth level needs coherent sources}
\centering
\card{height=5.6cm}{step5_joint.pdf}
\take{Best $r$ with mid within 0.02 of 0.848: 0.845 at mid 0.828 (edge; coherent). Independent: mid $\le$ 0.665.}
\src{Numerical optimum; maximality unproven; band-mean level 0.946 out of reach; at $\sigma=0.20$ both meet it (mid 0.868 coherent, 0.831 independent).}
\end{frame}

% 18
\begin{frame}{Step 6: held-out areas reach $r\ge0.6$: coherent 11/11, independent 7/11}
\centering
\card{height=5.6cm}{step6_sigma_loao.pdf}
\take{$\sigma=0.15$, $n=11$ folds, mean $\pm$ SD: coherent $0.81\pm0.05$, independent $0.62\pm0.09$.}
\src{10-area mean vs held-out: $0.93\pm0.06$, above the fitted readout's 0.81. Tests the target fit only. Folds share 9 of 10 training areas; SD descriptive.}
\end{frame}

% 19
\begin{frame}{Step 7: the circuit is far from the readout optimum}
\centering
\card{height=5.55cm}{step7_circuit_vs_opt.pdf}
\take{Circuit: S 63--78\%, M 17--19\% of gamma power, Re coherence $-0.06$ to 0.27; optimum 46/49/5\%, $\pm1$. At the declared depths and under this readout, the gap is in the circuit's spectra.}
\src{ref s1--s3: reference column, seeds 1--3; selected s4--s6: selected setting, fresh seeds 4--6. 50--150 Hz unit cross-spectra. Conditional on the three-source readout; its unit-dependent filtering is unvalidated.}
\end{frame}

% geometry
\begin{frame}{Readout depth is a confound: moving it alone nearly doubles map similarity}
\centering
\card{height=5.0cm}{geometry.pdf}
\take{Spectra fixed: map $r$ 0.38$\to$0.72, not distinguished from a frequency-shift null ($p\approx0.19$, 20 draws, best of 1330 each).}
\src{Depths S $-0.25$, M $-0.05$, D 0.20. Null on the incoherent spectra (unshifted best 0.725, the baseline here): unit spectra circularly shifted, power kept; 12/20 met none, 3 beat 0.725; 20 draws is a weak test. At the selected depths $r_{ab}$ reached 0.6 in 1 of the 8 other runs.}
\end{frame}

% 20
\begin{frame}{The tested density-informed column does not reproduce the primate motif}
\vspace{0.3cm}
\centering
\begin{minipage}{0.95\textwidth}\raggedright\small
\begin{itemize}\setlength{\itemsep}{0.6em}
\item Made: alpha-beta in the deep half on the reference seeds ($r_{ab}$ 0.83--0.86), but flat density does too (0.84). It rests on the imposed SOM/PV ramp, which runs against the measured CB/PV profile; reversing it breaks it. The selected setting lost it on fresh seeds (0.35--0.47); the search selected on whole-map similarity, not the band criteria. The reference was not rerun.
\item Disrupted: changing density or the ramp hands the map to one unit; no band moves alone.
\item Missing: gamma near L4, 0.19--0.31 vs 0.848. Weight changes reached $r_{gm}\ge0.6$ only without the crossing; with it, at most 0.54. Every setting fails the scorer.
\item Readout ($\sigma=0.15$, declared depths): coherent sources reach $r$ 0.853 (0.845 with mid near target), independent 0.657; the circuit's Re coherence is $-0.06$ to 0.27. Moving the depths alone: map $r$ 0.38$\to$0.72, not distinguished from a shift null.
\end{itemize}
\end{minipage}
\vfill
\end{frame}

\begin{frame}{Scope and limits}
\vspace{0.3cm}
\centering
\begin{minipage}{0.95\textwidth}\raggedright\small
\begin{itemize}\setlength{\itemsep}{0.6em}
\item Tested: one three-unit E/PV/SOM column. Varied: unit sizes, in-degree rule, imposed SOM/PV ramp, SOM$\to$E kinetics, network size. Not tested: PV$\to$E kinetics, E/I balance, PV or SOM silencing, VIP cells.
\item The model is stationary: every unit gets the same stationary Poisson drive, while the observed maps come from timed peri-stimulus inputs.
\item Target: the LANSA mean; the motif's generality is disputed.
\item The readout conclusion is conditional on the three-source readout. This says nothing about where gamma is generated in macaque cortex.
\item Next: a middle-layer gamma generator (PV kinetics, rate-matched control) and source geometry.
\end{itemize}
\end{minipage}
\vfill
\end{frame}

% 21
\begin{frame}{Sources}
\fontsize{8}{9.5}\selectfont
\begin{tabular}{@{}lp{0.82\textwidth}@{}}
\toprule
Page & Manuscript item; records (Table S6) and data \\
\midrule
2--3 & Fig.~1; LANSA densities and LFP (Mulvey et al.\ 2024); builder C1 \\
4 & Table~1; isolated-unit check R1, seal S1; SOM/PV sweep R2, seal S2; unit peaks R4 \\
5--7 & Fig.~2, Results; run rates R3; size-ladder verdicts R6; mid-depth gamma D1 (code C3); cross terms D2 (C4) \\
8--10 & Fig.~3, Fig.~S4; scorer C2; ceilings R5; verdicts R6, R7, R8; search seal S4, selections R3 \\
11 & Fig.~5, Table~5; exploratory weight runs R8 \\
12--13 & Fig.~6, Table~6; controls seal S7; per-run outcomes D8 (C12); engine check D9; maps R11 \\
15--19 & Results, readout capacity; D3 (C5); joint constraint D5 (C7); grid check D7 (C11) \\
20 & Table~8; leave-one-area-out D4 (C6) \\
21 & Table~7; circuit vs optimum D5 (C7) \\
22 & Results, geometry; findings log N1 \S7 \\
\bottomrule
\end{tabular}
\vfill
\src{Manuscript in preparation (draft 2026-11-10). Figures redrawn from these records by the deck's figure script.}
\end{frame}

\end{document}
